\begin{table*}[t]
\centering
\caption{Status of HEA biofilm/MIC evidence and a proposed program for AM HEAs. Existing studies are on cast/wrought material in simulated media; no AM-HEA natural-seawater biofilm study exists yet.}
\label{tab:biofilm}
\footnotesize
\begin{tabularx}{\textwidth}{@{}lX@{}}
\toprule
Element & Design \\
\midrule
Existing evidence & \emph{P.~aeruginosa} on FeCoCrNiMo$_{0.1}$: passive film thinned, selective Fe dissolution via EET (first HEA MIC study) \cite{lou2020mic}; on Cantor-type HEA: max pit 4.77 vs 3.16~$\mu$m (14~d), OCP $-367\!\to\!\sim\!-500$ mV \cite{mic_cantor2022}; \emph{D.~vulgaris} (SRB) on CoCrFeMnNi: sulfides, film thinning, more aggressive under carbon starvation \cite{wang2023srb} \\
Materials & LPBF CoCrFeMnNi, SLM AlCoCrFeNi$_{2.1}$, and one WAAM deposit; as-built vs machined; LPBF 316L and 2205 as controls \\
Organisms/media & \emph{P.~aeruginosa} (aerobic model), \emph{D.~vulgaris} (SRB), natural-seawater biofilm controls; 2216E and ASTM D1141 media \cite{astm_d1141}; carbon-starved variants \\
Electrochemistry & OCP, LPR and EIS time series over 14--28 days, inoculated vs sterile (per \cite{mic_cantor2022,wang2023srb} protocols); $n\geq3$ \\
Surface/film analysis & CLSM live/dead imaging, SEM pit mapping, XPS/AES depth profiles of passive film, SKPFM micro-galvanic maps \\
Under-film chemistry & pH and O$_2$ microelectrodes beneath biofilms at defined sites; sulfide assays for SRB cells \\
Integration & Couple one coupon set into the natural-seawater racks of Table~\ref{tab:exposure} so laboratory MIC and field biofilm data converge \\
\bottomrule
\end{tabularx}
\end{table*}
